\section{Evidence-Gated Methodology}\label{sec:method}

This study is as much about method as about results. Every experiment was
executed by an autonomous research loop on a single machine, governed by
machine-checkable contracts that decide what may launch, what may be measured,
and what a number is allowed to be called. The notation and verdict semantics
defined here are used by all results sections.

\subsection{The Unit of Claim}

The atomic claim of this paper is a \emph{paired across-seed bits-per-byte
(BPB) delta with a 95\% confidence interval, reported on at least two corpora,
under a single-variable diff at iso-FLOP, scored by a versioned harness against
a measured noise floor}. Each ingredient is defined below.

\textbf{Metric.} For a model $m$ and a fixed evaluation corpus $C$ of byte
length $B(C)$, tokenized by $m$'s own tokenizer into $x_1,\dots,x_N$,
\begin{equation}
\mathrm{BPB}(m,C) \;=\; \frac{1}{B(C)} \sum_{t=1}^{N} -\log_2 p_m\!\left(x_t \mid x_{<t}\right).
\end{equation}
Because the denominator is a property of the raw text rather than of any
tokenization, BPB remains comparable when tokenizers differ and across corpora,
unlike per-token perplexity. The harness scores fixed windows of 204{,}600
tokens per corpus over 869{,}710 bytes of WikiText-2 and 843{,}643 bytes of
held-out Python code, at pinned dataset revisions, with 13-gram decontamination
against training data; per-checkpoint perplexities and noise floors on these
corpora appear in Table~\ref{tab:arms}. Effects must be reported on both
corpora: $\Delta_c = \mathrm{BPB}_{\mathrm{base}}(c) -
\mathrm{BPB}_{\mathrm{treat}}(c)$, positive meaning improvement.

\textbf{Single variable at iso-FLOP.} Two arms compare only if exactly one
variable differs \emph{and} their training FLOPs, under the
$6N + 12 \cdot L \cdot H \cdot Q \cdot T$ accounting of \citet{chowdhery2022palm},
match within 5\%; comparisons at unequal compute conflate the variable with
budget \citep{hoffmann2022training}. Token-matched is not FLOP-matched when the
architecture or optimizer cost changes, so the check is computed, not assumed,
and recorded with each run (the largest recorded mismatch among the pretraining
comparisons is a FLOPs-per-token ratio of 1.00043; Table~\ref{tab:attrib}).
Every comparison states its matched budget in tokens and steps
(Table~\ref{tab:hparams}).

\textbf{Paired seed test.} Both arms train the identical seed set
$\{0,1,2\}$, varying initialization and data order. The effect is the
difference of arm means with a Student-$t$ 95\% interval,
$\Delta_c \pm t_{0.975,\nu}\sqrt{\mathrm{SEM}_b^2+\mathrm{SEM}_t^2}$, where
$\nu$ is the Welch--Satterthwaite degrees of freedom floored to an integer
(conservative at $n{=}3$, where a normal quantile would be too narrow). A delta
is \emph{significant} only if the entire interval lies on the improving side of
zero. A corpus-subsample noise floor -- the spread of the baseline checkpoint
over three disjoint evaluation subsamples -- is retained as a secondary check;
deltas below it are labeled not significant. The floors themselves motivate the
design: on this suite the code-corpus perplexity floor is large
(Table~\ref{tab:arms}), so single-checkpoint perplexity gaps on code are close
to meaningless, whereas across-seed BPB intervals are not. Results are quoted
as $+X$ BPB (95\% CI $[\mathrm{lo},\mathrm{hi}]$, 3 seeds, paired).

\textbf{Versioned harness.} Cross-run comparable numbers come only from the
evaluation harness, which stamps a suite version into every result:
(text-lm-v2) for the BPB/perplexity suite, (math-acc-v1) for the pinned math
answer-extractor and verifier, (ecl-ladder-v1) for the passkey context ladder.
Any suite change bumps the version, so numbers with different stamps are never
differenced. Perplexities printed by the training loop are labeled in-loop
validation perplexity, single run, and are advisory only. This distinction is
load-bearing: an in-loop 0.68-perplexity SFT ``win'' collapsed to $+0.009$
(95\% CI $[0.0045, 0.0139]$) once re-scored on one fixed held-out token set
(Figure~\ref{fig:sft-confound}).

\subsection{Claim Classes and Verdict Caps}

Verdicts take the values \emph{win}, \emph{loss}, \emph{inconclusive}, and
\emph{directional}. Every run carries a lifecycle stage (data, architecture,
mid-training, sft, rlvr, \dots), and a per-stage registry names the required
(HARD) evaluation battery for that stage. A missing HARD item forces the
verdict to directional regardless of the point estimate -- completeness gates
significance, not the reverse. A single-seed run is never significance-testable
and is likewise capped at directional; the RLVR cohort was run at one seed by
pre-registered design and carries exactly this cap (Table~\ref{tab:rlvr}).
Near-chance downstream results are reported with Wilson intervals as no-signal,
never as wins, with pass@$k$ computed by the unbiased estimator of
\citet{chen2021evaluating}.

\subsection{Control Reuse and Pre-Registration}

Baseline checkpoints are reused across experiments as symlinks and re-scored
in-cohort, so every number in a comparison comes from the same evaluation pass
of the same harness version: the architecture sub-drill reuses the de-confound
baselines (Table~\ref{tab:attrib}), and the data-mix study reuses the dclm-edu
arm (Table~\ref{tab:data}). Expected outcomes are written down before launch
and bind in both directions. The RLVR plan, written 2026-07-01 before the
2026-07-02 launch, predicted a null at this scale
\citep{yue2025does} and mandated a random-reward control for Qwen-lineage
models \citep{shao2025spurious}; the null was confirmed
(Figure~\ref{fig:rlvr}). The mid-training anneal pre-registered an expected
English null that the data violated -- $+0.0157$ BPB (95\% CI
$[+0.0140, +0.0174]$, 3 seeds, paired; Table~\ref{tab:midtrain}) -- and the
violation is reported as exploratory rather than promoted to a claim.

\subsection{Launch Evidence and the Ledger}

No budgeted GPU run starts until its evidence is recorded: a literature-sourced
token budget, the arm~$\times$~seed plan, the confound check, a measured
throughput probe with a derived ETA, and a passing smoke test that executes the
exact training script for one step on a tiny batch and asserts a finite,
descending loss, a checkpoint save-then-reload round trip, and exit code zero.
The ledger is the append-only truth store: a single programmatic interface with
atomic writes records techniques, runs, verdicts, suite stamps, and costs, and
is the deduplication authority for what has already been tried. A resumable
state machine executes the stages unattended -- preflight, scan, selection,
briefing, data preparation, launch, evaluation, verdict -- and always ends by
writing a digest, while irreversible actions (new model builds, compute spend,
publication) remain human-gated.

\subsection{Safety Envelope}

All training ran on one NVIDIA GB10 with a unified CPU+GPU memory pool of
$\approx$119~GB and no separate VRAM; over-allocation invokes the kernel
OOM-killer and can take down the machine, so safety limits are part of the
method. Three constraints apply to every job: (i) a per-process allocator guard
caps GPU memory at a 0.85 fraction of the pool and is imported before the
framework in every script; (ii) an independent watchdog process kills any
unattended trainer at a 0.80 pool fraction -- strictly below the guard, so the
watchdog fires first -- and only one trainer may run at a time; (iii)
cross-entropy over the 151{,}936-token vocabulary is computed in chunked fp32,
never materializing the full sequence-by-vocabulary logit matrix. Throughput
and MFU figures are reported as approximate throughout, because the device peak
for this hardware is an estimate.
